% =====================================================================
% ECN-BENCH — Data Availability section
%
% STATUS: ALREADY APPLIED. As of 2026-08-26 this section is live in
%   08_paper/ecn_bench_paper.tex, and the macros below are in its preamble.
%   The paper compiles clean at 87 pages, 0 errors, 0 undefined references,
%   0 overfull boxes. Only \ecnarchivedoi is still a placeholder.
%
%   This standalone copy is kept for reference, and for pasting into a
%   different manuscript. Editing it does NOT change the paper.
%
% HOW TO USE
%   1. Upload ECNBENCH-RELEASE/ to Zenodo, OSF, Figshare or a Dataverse.
%      Zenodo gives a permanent, versioned DOI and accepts up to 50 GB;
%      this archive is ~2.4 GB.
%   2. Fill in \ecnarchivedoi and \ecnarchiveurl below.
%   3. Replace the existing \section*{Data Availability} in
%      ecn_bench_paper.tex with everything after the RULE line.
%
% ANONYMOUS REVIEW: leave \ecnanonymoustrue set. The section then routes
% reviewers through the venue and defers the public link to camera-ready,
% exactly as the current draft does. Flip it to \ecnanonymousfalse for the
% camera-ready and the real link appears everywhere it is needed.
% =====================================================================

\newif\ifecnanonymous
\ecnanonymoustrue          % <-- \ecnanonymousfalse for camera-ready

\newcommand{\ecnarchivedoi}{10.5281/zenodo.XXXXXXX}    % <-- your DOI
\newcommand{\ecnarchiveurl}{https://doi.org/\ecnarchivedoi}
\newcommand{\ecnarchive}{%
  \ifecnanonymous
    the replication archive accompanying this submission%
  \else
    the replication archive (\href{\ecnarchiveurl}{\texttt{\ecnarchivedoi}})%
  \fi
}

% Requires \usepackage{hyperref} and \usepackage{url} in the preamble.

% ===================== RULE — copy from here =========================

\section*{Data Availability}

Every quantitative claim in this paper is derived from artifacts released in
\ecnarchive. The archive is arranged in reading order and each directory carries its own
\texttt{README}; \texttt{00\_docs/PROVENANCE.md} records which claims survived the audit
of Section~\ref{sec:evaluator_replication} and which were withdrawn, and
\texttt{00\_docs/REPRODUCE.md} gives the exact commands that regenerate every table and
figure below. Those commands were verified end to end against the archive: all five
analyses reproduce byte-identical output.

\textbf{Campaign records.} For each of the four campaigns we release the per-unit scored
records: \texttt{event\_results.json} (the July 2026 pass, Evaluator~J,
DeepSeek-R1-14B), \texttt{event\_results\_evaluatorB\_openrouter.json} together with
\texttt{event\_results\_evaluatorA\_repeat\{1,2,3\}.json} (the four independent
Evaluator~A reads of Section~\ref{sec:four_model_extension}), and
\texttt{event\_results\_evaluatorC\_gptlunapro.json} (the independent Evaluator~C read).
The passes are kept in separate files rather than merged, because the derived metrics
stored in each are valid only for that pass's own probability vectors. The rollup behind
Section~\ref{sec:four_model_extension} and Table~\ref{tab:four_model_clean} is
\texttt{reproducibility\_analysis\_master.json}, which is the source of truth for those
numbers.

\textbf{Simulation transcripts.} The run-integrity audits of
Appendices~\ref{sec:persona_audit}--\ref{sec:evidence_depth} read the per-unit SQLite
trace databases. These are released in full for Llama-3.1-8B, Mistral-7B and
Qwen2.5-7B ($120$ databases per model). For Qwen2.5-14B they are released for the six
events re-simulated 2026-08-21 ($18$ units), together with a three-unit fragment of the
original 2026-07-27 campaign (event~C15); the transcripts for that campaign's remaining
$21$ events were removed by the cluster's scratch-filesystem retention policy before
they could be transferred, and are unrecoverable. Their scored records are unaffected
and are released in full. One consequence is documented rather than repaired: the
all-campaign action-density counts of Table~\ref{tab:action_density_30event} come from a
2026-07-27 cluster-side scan of Qwen2.5-14B's then-30-event campaign, six of which were
later found to have been simulated under the wrong topic
(\S\ref{sec:fourth_model_status}); that scan could not be redone against the corrected
traces.

\begin{sloppypar}
\textbf{Post-campaign experiments.} Each is released with its generator, its analysis and
its figure script, so that every number can be recomputed from the raw runs rather than
taken from a summary. The cross-evaluator degeneracy-transfer probe of
Section~\ref{sec:volume_mechanism} is \texttt{run\_\allowbreak degeneracy\_\allowbreak
transfer.py} and \texttt{analyze\_\allowbreak degeneracy\_\allowbreak transfer.py} over
\texttt{degeneracy\_\allowbreak transfer.json}. The $2\times4$ truncation grid is
\texttt{build\_\allowbreak truncation\_\allowbreak inputs.py},
\texttt{run\_\allowbreak truncation\_\allowbreak grid.py} and
\texttt{analyze\_\allowbreak truncation\_\allowbreak grid.py} over
\texttt{event\_\allowbreak results\_\allowbreak trunc\_*.json} and
\texttt{truncation\_\allowbreak results.json}. The single-agent baseline of
Section~\ref{sec:cot_baseline} is \texttt{run\_\allowbreak single\_\allowbreak
agent\_\allowbreak cot\_\allowbreak baseline.py}, producing
\texttt{single\_\allowbreak agent\_\allowbreak cot\_\allowbreak baseline\_\allowbreak
results.json} ($240$ runs, all fields retained except the chain-of-thought text, which
is truncated to $300$ characters per run), analysed by
\texttt{analyze\_\allowbreak cot\_\allowbreak baseline.py} into
\texttt{cot\_\allowbreak baseline\_\allowbreak analysis.json}. That last file carries all
three CoT aggregation variants and both pooling scopes, not only the ones tabulated
here, so a reader can check the sensitivity of \S\ref{sec:cot_baseline} directly.
\end{sloppypar}

\textbf{Benchmark and pipeline.} The event catalogue, seed dossiers, injection bank and
scoring rubric are released so the benchmark can be run against new models, alongside
the simulation engine, the HPC protocol runner and the SLURM job scripts as executed.
The archive also retains the superseded and invalid artifacts the audit produced --- the
defective repair, the wrong-topic re-simulation, the two divergent event catalogues ---
so that every defect reported in \S\ref{sec:outage} and \S\ref{sec:fourth_model_status}
is independently checkable rather than merely asserted.

\ifecnanonymous
To preserve anonymity during review, the archive is not linked here; it is available to
reviewers through the venue, and the permanent public record will be linked in the
camera-ready version.
\else
The archive is permanently available at \url{\ecnarchiveurl}.
\fi

% ===================== to here ======================================


% ---------------------------------------------------------------------
% OPTIONAL: a footnote for the first page, if the venue allows one.
% ---------------------------------------------------------------------
% \thanks{Data, code and audit trail: \ecnarchive.}

% ---------------------------------------------------------------------
% OPTIONAL: a bibliography entry, if you prefer to cite the archive.
% ---------------------------------------------------------------------
% \bibitem{ecnbench-archive}
% ECN-BENCH replication archive: campaign records, simulation transcripts,
% post-campaign experiments and audit trail.
% \newblock Zenodo, 2026. \newblock DOI: \ecnarchivedoi.
